\honest{Belief in the Implied Invisible}{Eliezer Yudkowsky}{2008}

The anti-zombie argument does not teach that ``Anything you can't see doesn't exist.'' Suppose I send a photon into a void between superclusters. Because the expansion of the universe is accelerating, a time will come after which nothing from it can ever reach me. Even aliens who caught it could not come back fast enough. Does it vanish then? No. That would break conservation of energy, the second law and just about every other law, and would imply the photon knows exactly when I stop being able to see it. It is ``a silly idea.'' \nb{The physics is right. In the standard cosmological model there is an event horizon, and the photon crosses it after a finite time.} But then why may I not believe in any invisible thing I like?

I remember reading, though I could not find a source, that Occam's razor was once invoked against the idea that the Milky Way is made of millions of stars, or perhaps that the nebulae are galaxies full of stars, because it multiplied entities. \nb{The editor could not find one either.} I quote the razor in Latin, ``Entia non sunt multiplicanda praeter necessitatem,'' and call it ``Occam's original formulation.'' \nb{That phrasing is not found in Ockham's surviving works; it comes from John Punch, in 1639. Ockham's own versions say much the same.}

Extra stars do not count against a theory. Neither Bayesian formalization of the razor, Solomonoff induction or Minimum Message Length, penalizes galaxies for being big, and they had better not: reality keeps turning out bigger, from the Earth at the centre to Avogadro's number, and a razor that always counted against size would have been consistently wrong. In Solomonoff induction, what counts is the length of the program, not the memory or the running time it needs, and ``Occam's Razor doesn't care about RAM.'' \nb{That holds for the two formalizations the post names. Schmidhuber's speed prior, published in 2002, does count running time.}

In Minimum Message Length, the message states the equations, not the position of every quark. If the equations take 100 bits, there are about $2^{100}$ models of that size, so you take a $2^{-100}$ prior penalty and need about 100 bits of evidence, and more galaxies add nothing, unless your predictions depend on exact initial conditions, in which case the extra quarks do count. A vanishing photon, by contrast, needs an extra law, and ``It's the laws that are `entities'.''

Is the answer, then, just that my priors favour the photon's survival? That is ``not quite right.'' We assign probability to simple, well-tested equations and believe what they logically imply. That is belief in the ``implied invisible.'' Believing that the Flying Spaghetti Monster ate the photon, even just this once, or without reason that it hit a dust speck, is belief in a specific extra event. Believing such things happen in general is belief in an extra law. That is belief in the ``additional invisible,'' and gets no more than its unaltered prior. If a belief is something you track and count evidence for, perhaps you should hold no such beliefs at all. Ruling out everything we cannot interact with would be simpler, but very silly. \nb{The extra law is still rejected for costing prior probability, so the new account moves the simplicity answer onto the laws rather than replacing it. Max Tegmark had made the central point in 2003: unobservable parts do not make a theory untestable.}

This rule ``does seem'' to rule out all abuses. There are exotic cases that break it in theory, such as epiphenomenal demons who will torture $3\uparrow\uparrow\uparrow 3$ victims somewhere you can never verify if you say ``Niblick,'' but I cannot think of a case where it fails ``in human practice.''

In an addendum I ask whether you would send colonists at nearly light speed to a supercluster from which, by the time they arrive, no message could come back. Is the purely altruistic effort worth it, for the people who will live there happily, or does the ship blip out of existence on the way? This could become a very real question. \nb{Yudkowsky later uses the principle to argue for many worlds in quantum mechanics, in ``Decoherence is Simple.''}
